% Curriculum Vitae — Yoichi Ishibashi
% Build: lualatex cv.tex  (run twice for page numbers/links)
% Then copy cv.pdf to ../assets/pdf/cv.pdf and push to deploy.
\documentclass[11pt,letterpaper]{article}
\usepackage{amsmath} % provides \smash[b]
\usepackage{luatexja}
\usepackage[T1]{fontenc}
\usepackage{tgschola}
\usepackage{geometry}
\geometry{letterpaper, left=2.5in, right=1.05in, top=0.9in, bottom=1.1in}
\usepackage{fancyhdr}
\usepackage[normalem]{ulem}
\usepackage[colorlinks=true, urlcolor=blue, linkcolor=blue]{hyperref}
\urlstyle{same}

% Underline links like the Google Docs original (rule inherits link color)
\let\oldhref\href
\renewcommand{\href}[2]{\oldhref{#1}{\uline{#2}}}
\renewcommand{\url}[1]{\oldhref{#1}{\uline{#1}}}

\pagestyle{fancy}
\fancyhf{}
\renewcommand{\headrulewidth}{0pt}
\fancyfoot[R]{\scriptsize Curriculum Vitae, Yoichi Ishibashi, \thepage}

\setlength{\parindent}{0pt}
% Left-aligned text without hyphenation, like the original
\AtBeginDocument{\raggedright}

\newlength{\labelcol}\setlength{\labelcol}{1.2in}
\newlength{\labelgap}\setlength{\labelgap}{0.21in}

% Section: horizontal rule + label hanging in the left margin
% (\smash[b] keeps a multi-line label from pushing the next line down)
\newcommand{\cvsection}[1]{%
  \par\vspace{26pt}%
  {\noindent\rule{\textwidth}{2.4pt}\par}\nopagebreak\vspace{9pt}%
  \noindent\llap{\smash[b]{\parbox[t]{\labelcol}{\raggedright\large\bfseries #1}}\hspace{\labelgap}}\ignorespaces}

% Bold entry heading (\entryfirst: first entry right after \cvsection, no leading gap)
\newcommand{\entry}[1]{\par\vspace{10pt}\noindent\textbf{#1}\par\nopagebreak}
\newcommand{\entryfirst}[1]{\textbf{#1}\par\nopagebreak}

% Indented detail block under an entry
\newenvironment{indented}{%
  \par\nopagebreak
  \begin{list}{}{\setlength{\leftmargin}{0.3in}\setlength{\topsep}{1pt}%
    \setlength{\partopsep}{0pt}\setlength{\itemsep}{0pt}\setlength{\parsep}{1pt}}\item[]}{\end{list}}

% Right-aligned line (awards)
\newcommand{\awardline}[1]{{\par\raggedleft #1\par}}

% Publication entry: authors / title / venue (\paperfirst: first one after \cvsection)
\newcommand{\paper}[3]{\par\vspace{12pt}\noindent #1\\ #2\\ {\itshape #3}\par}
\newcommand{\paperfirst}[3]{#1\\ #2\\ {\itshape #3}\par}

\begin{document}

% ---------- Header ----------
\noindent\hspace*{-1.5in}\parbox{\dimexpr\textwidth+1.5in\relax}{%
  {\huge\bfseries Yoichi Ishibashi, Ph.D.}\hfill
  {\textbf{Last Updated:} Jun.~12, 2026}}

% ---------- Biography ----------
\cvsection{Biography}%
Yoichi Ishibashi is a researcher at NEC.\\
He received a bachelor's degree from Kyoto Sangyo University, Kyoto, Japan,
in 2018, and master's and Ph.D.\ degrees from Nara Institute of Science and
Technology (NAIST) in 2020 and 2023, respectively.

\vspace{8pt}
His research interest lies in Coding Agent and Self-Improving LLM.

% ---------- Contact ----------
\cvsection{Contact Information}%
\textbf{Affiliation:} Generative AI Group,\\
\hspace*{0.55in}\href{https://www.nec.com/en/global/rd/}{Knowledge Science Research Laboratories},\\
\hspace*{0.55in}\href{https://www.nec.com/en/global/rd/}{NEC Corporation}\\
\textbf{Email:} yoichi-ishibashi@nec.com\\
\textbf{Website:} \url{https://yoichi1484.github.io/}

% ---------- Education ----------
\cvsection{Education}%
\entryfirst{Nara Institute of Science and Technology}
\begin{indented}
\href{https://ahcweb01.naist.jp/en/}{Augmented Human Communication Laboratory}\\
April 2020 - March 2023, Nara, Japan\\[3pt]
Doctor of Engineering\\[3pt]
Thesis title: ``Semantic Operations on an Embedding Space''\\[3pt]
{\itshape Supervisor: Professor \href{https://ahcweb01.naist.jp/Prof.Nakamura/index_e.html}{Satoshi Nakamura}}\\
{\itshape Subadvisor: Associate Professor \href{https://www.sudoh.nl/}{Katsuhito Sudoh}}
\end{indented}

\entry{Nara Institute of Science and Technology}
\begin{indented}
\href{https://ahcweb01.naist.jp/en/}{Augmented Human Communication Laboratory}\\
April 2018 - March 2020, Nara, Japan\\[3pt]
Master of Engineering\\[3pt]
Thesis title: ``Reflection-based Word Attribute Transfer''\\[3pt]
{\itshape Supervisor: Professor \href{https://ahcweb01.naist.jp/Prof.Nakamura/index_e.html}{Satoshi Nakamura}}\\
{\itshape Subadvisor: Associate Professor \href{https://www.sudoh.nl/}{Katsuhito Sudoh}}
\end{indented}

\entry{Kyoto Sangyo University}
\begin{indented}
April 2014 - March 2018, Kyoto, Japan\\[3pt]
Bachelor of Computer Science and Engineering\\[3pt]
Thesis title: ``Generating Responses based on Information Visually-Induced by Text Utterance''\\[3pt]
{\itshape Supervisor: Professor \href{https://researchmap.jp/read0006204?lang=en}{Hisashi Miyamori}}\\
President's Award (\textbf{1st in Class})
\end{indented}

% ---------- Research Interests ----------
\cvsection{Research Interests}%
My research interests include \textbf{coding agents} and \textbf{self-improving LLMs}.

% ---------- Research Experience ----------
\cvsection{Research Experience}%
\entryfirst{Researcher}
\begin{indented}
NEC Corporation, Generative AI Group,\\
Knowledge Science Research Laboratories, Japan\\
April 2024 - Present\\
Research Topic: \href{https://arxiv.org/abs/2410.15639}{Self-improving LLM}, \href{https://arxiv.org/abs/2505.10182}{Reasoning model}
\end{indented}

\entry{Program-Specific Researcher}
\begin{indented}
Kyoto University, Kyoto, Japan\\
April 2023 - March 2024\\
Research Topic: \href{https://arxiv.org/abs/2309.11852}{Knowledge Editing} and \href{https://arxiv.org/abs/2404.02183}{LLM Agent}
\end{indented}

\entry{Research Intern (Visiting Research Student)}
\begin{indented}
Advisor: {\itshape Prof.\ \href{https://danushka.net/}{Danushka Bollegala}}\\
The University of Liverpool, Department of Computer Science,\\
Natural Language Processing Group, Liverpool, UK\\
February 2022 - July 2022\\
Research Topic: \href{https://arxiv.org/abs/2302.05619}{Prompt Learning}
\end{indented}

\entry{Research Assistant}
\begin{indented}
with {\itshape Prof.\ \href{https://ahcweb01.naist.jp/Prof.Nakamura/index_e.html}{Satoshi Nakamura}}\\
Nara Institute of Science and Technology, Nara, Japan\\
April 2020 - March 2023
\end{indented}

\entry{Researcher}
\begin{indented}
Nara Institute of Science and Technology, Nara, Japan\\
October 2019 - March 2023\\
Research Topic: \href{https://arxiv.org/abs/2210.13034}{Word Embeddings}
\end{indented}

% ---------- Teaching ----------
\cvsection{Teaching Experience}%
\textbf{Teaching Assistant} at the Nara Institute of Science and Technology
\begin{indented}
Sequential Data Modeling, Fall 2019\\
Data Engineering, Fall 2019\\
Machine Learning, Spring 2019
\end{indented}

\vspace{6pt}
\textbf{Teaching Assistant} at the Kyoto Sangyo University
\begin{indented}
Multimedia Coding, Fall 2017
\end{indented}

% ---------- Research Papers ----------
\cvsection{Research Papers}%
\paperfirst{\textbf{Yoichi Ishibashi}, Taro Yano, Masafumi Oyamada.}
{\href{https://arxiv.org/abs/2605.15221}{Effective Harness Engineering for Algorithm Discovery with Coding Agents}}
{In ICML 2026 AI for Science Workshop (\textbf{ICML 2026 Workshop}), 2026.}

\paper{Haowen Li, \textbf{Yoichi Ishibashi}, Masafumi Oyamada.}
{Evaluating the Impact of Reviewer Guideline Design on LLM-Based Automated Peer Review}
{In Findings of the Association for Computational Linguistics (\textbf{ACL 2026 Findings}), 2026.}

\paper{Taro Yano, \textbf{Yoichi Ishibashi}, Masafumi Oyamada.}
{\href{https://arxiv.org/abs/2505.21963}{LaMDAgent: An Autonomous Framework for Post-Training Pipeline Optimization via LLM Agents}}
{In Proceedings of the 2025 Conference on Empirical Methods in Natural Language Processing (\textbf{EMNLP 2025}), November 2025.}

\paper{\textbf{Yoichi Ishibashi}, Taro Yano, Masafumi Oyamada.}
{\href{https://arxiv.org/abs/2505.10182}{Mining Hidden Thoughts from Texts: Evaluating Continual Pretraining with Synthetic Data for LLM Reasoning}}
{Preprint, May.\ 2025.}

\paper{Yunzhen He, Yusuke Takase, \textbf{Yoichi Ishibashi}, Hidetoshi Shimodaira.}
{\href{https://arxiv.org/abs/2503.02343}{DeLTa: A Decoding Strategy based on Logit Trajectory Prediction Improves Factuality and Reasoning Ability}}
{Proceedings of the 2nd Workshop on Uncertainty-Aware NLP (\textbf{UncertaiNLP 2025}), March.\ 2025.}

\paper{\textbf{Yoichi Ishibashi}, Taro Yano, Masafumi Oyamada.}
{\href{https://arxiv.org/abs/2410.15639}{Can Large Language Models Invent Algorithms to Improve Themselves?}}
{In Proceedings of the 2025 Conference of the Nations of the Americas Chapter of the Association for Computational Linguistics (\textbf{NAACL 2025}), March 2025.}

\paper{\textbf{Yoichi Ishibashi}, Yoshimasa Nishimura.}
{\href{https://arxiv.org/abs/2404.02183}{Self-Organized Agents: A LLM Multi-Agent Framework toward Ultra Large-Scale Code Generation and Optimization}}
{Preprint, March.\ 2024.}

\paper{\textbf{Yoichi Ishibashi}, Sho Yokoi, Katsuhito Sudoh, Satoshi Nakamura.}
{\href{https://arxiv.org/abs/2210.13034}{Subspace Representations for Soft Set Operations and Sentence Similarities}}
{In Proceedings of the 2024 Conference of the North American Chapter of the Association for Computational Linguistics (\textbf{NAACL 2024}), June 2024.}

\paper{\textbf{Yoichi Ishibashi}, Hidetoshi Shimodaira.}
{\href{https://arxiv.org/abs/2309.11852}{Knowledge Sanitization of Large Language Models}}
{Preprint, September.\ 2023.}

\paper{\textbf{Yoichi Ishibashi}, Danushka Bollegala, Katsuhito Sudoh, Satoshi Nakamura.}
{\href{https://arxiv.org/abs/2302.05619}{Evaluating the Robustness of Discrete Prompts}}
{In Proceedings of the 17th Conference of the European Chapter of the Association for Computational Linguistics (\textbf{EACL 2023}), May 2023.}

\paper{Kosuke Takahashi, \textbf{Yoichi Ishibashi}, Katsuhito Sudoh, Satoshi Nakamura.}
{\href{http://www.statmt.org/wmt21/pdf/2021.wmt-1.113.pdf}{Multilingual Machine Translation Evaluation Metrics Fine-tuned on Pseudo-Negative Examples for WMT 2021 Metrics Task}}
{In Proceedings of the Sixth Conference on Machine Translation (\textbf{WMT 2021}), 1031 - 1034, November 2021.}

\paper{\textbf{Yoichi Ishibashi}, Katsuhito Sudoh, Koichiro Yoshino, Satoshi Nakamura.}
{\href{https://www.jstage.jst.go.jp/article/jnlp/28/1/28_206/_article/-char/en}{Reflection-based Word Attribute Transfer.}}
{In Journal of Natural Language Processing, vol.28, no.1, March 2021.}

\paper{\textbf{Yoichi Ishibashi}, Katsuhito Sudoh, Koichiro Yoshino, Satoshi Nakamura.}
{\href{https://www.aclweb.org/anthology/2020.acl-srw.8/}{Reflection-based Word Attribute Transfer.}}
{In Proceedings of the 58th Annual Meeting of the Association for Computational Linguistics (\textbf{ACL 2020}): Student Research Workshop, 51 - 58, July 2020.}

\paper{\textbf{Yoichi Ishibashi}, Katsuhito Sudoh, Koichiro Yoshino, Satoshi Nakamura.}
{Reflection-based Word Attribute Transfer.}
{Third International Workshop on Symbolic-Neural Learning (\textbf{SNL 2019}), July 2019.}

\paper{\textbf{Yoichi Ishibashi}, Hisashi Miyamori.}
{\href{https://yoichi1484.github.io/projects/shizuka/doc/INLG2019/paper/paper_final2019.pdf}{Generating Responses based on Information Visually-Induced by Text Utterance.}}
{Technical Report, December 2017.}

\paper{\textbf{Yoichi Ishibashi}, Sho Sugimoto, Hisashi Miyamori.}
{KSU Teams Dialogue System at the NTCIR-13 Short Text Conversation Task 2.}
{In Proceedings of the 13th \textbf{NTCIR} Conference on Evaluation of Information Access Technologies, 248 - 256, December 2017.}

% ---------- Domestic Conferences ----------
\cvsection{Domestic Conferences}%
\paperfirst{何昀臻, 高瀬侑亮, 石橋陽一, 下平英寿}
{\href{https://www.anlp.jp/proceedings/annual_meeting/2024/pdf_dir/P6-21.pdf}{大規模言語モデルにおける幻覚緩和のための単語確率の外挿}}
{\upshape 言語処理学会 第30回 年次大会 (NLP 2024), リクルート賞 \textbf{(1/599)}}

\paper{石橋 陽一, 横井 祥, 須藤 克仁, 中村 哲}
{\href{https://www.anlp.jp/proceedings/annual_meeting/2023/pdf_dir/C4-1.pdf}{正準角および部分空間に基づく BERTScore の拡張}}
{\upshape 言語処理学会 第29回 年次大会 (NLP 2023)}

\paper{石橋 陽一, 横井 祥, 須藤 克仁, 中村 哲}
{\href{https://www.jstage.jst.go.jp/article/pjsai/JSAI2022/0/JSAI2022_1P4GS602/_article/-char/en}{学習済み埋め込み空間の線型部分空間を用いた集合演算}}
{\upshape 第36回 人工知能学会全国大会 (JSAI2022)}

\paper{石橋 陽一, 横井 祥, 須藤 克仁, 中村 哲}
{\href{https://www.anlp.jp/proceedings/annual_meeting/2022/pdf_dir/PH2-8.pdf}{線型部分空間に基づく学習済み単語埋込空間上の集合演算}}
{\upshape 言語処理学会 第28回 年次大会 (NLP 2022), 若手奨励賞 \textbf{(12/280)}}

\paper{高橋 洸丞, 石橋 陽一, 須藤 克仁, 中村 哲}
{\href{https://www.anlp.jp/proceedings/annual_meeting/2022/pdf_dir/G3-2.pdf}{単語属性変換で作成した疑似負例データを用いた自動機械翻訳評価}}
{\upshape 言語処理学会第28回年次大会 (NLP 2022)}

\paper{石橋 陽一, 横井 祥, 須藤 克仁, 中村 哲}
{学習済み埋め込み空間における集合演算}
{\upshape 第24回 情報論的学習理論ワークショップ (IBIS 2021)}

\paper{石橋 陽一, 横井 祥, 須藤 克仁, 中村 哲}
{量子論理に基づく単語埋込集合間の論理演算}
{\upshape NLP若手の会 第16回 シンポジウム (YANS 2021)}

\paper{石橋 陽一, 須藤 克仁, 中村 哲}
{\href{https://www.anlp.jp/proceedings/annual_meeting/2021/pdf_dir/P4-13.pdf}{単語属性変換による自然言語推論データの拡張}}
{\upshape 言語処理学会 第27回 年次大会 (NLP 2021)}

\paper{石橋 陽一, 須藤 克仁, 吉野 幸一郎, 中村 哲}
{\href{https://www.anlp.jp/proceedings/annual_meeting/2020/pdf_dir/P6-2.pdf}{鏡映変換に基づく埋め込み空間上の単語属性変換}}
{\upshape 言語処理学会 第26回 年次大会 (NLP 2020), 優秀賞 \textbf{(6/396)}}

\paper{石橋 陽一, 須藤 克仁, 吉野 幸一郎, 中村 哲}
{鏡映変換に基づく埋め込み空間上の単語属性変換}
{\upshape NLP若手の会 第13回 シンポジウム (YANS 2019), 奨励賞 \textbf{(7/93)}}

\paper{石橋 陽一, 須藤 克仁, 吉野 幸一郎, 中村 哲}
{鏡映変換に基づく埋め込み空間上の単語属性変換}
{\upshape 情報処理学会 第241回 自然言語処理研究会, 優秀研究賞 \textbf{(3/28)}}

\paper{石橋 陽一, 森 泰, 木村 輔, 宮森 恒}
{\href{https://www.anlp.jp/proceedings/annual_meeting/2019/pdf_dir/D3-4.pdf}{質問文から連想した画像特徴量を用いた質問応答モデル}}
{\upshape 言語処理学会 第25回 年次大会 (NLP 2019)}

\paper{杉本 翔, 石橋 陽一, 宮森 恒}
{\href{https://db-event.jpn.org/deim2018/data/papers/347.pdf}{ペルソナベクトルの演算に応じた新たな個性での対話応答文生成}}
{\upshape 第10回 データ工学と情報マネジメントに関するフォーラム (DEIM 2018)}

\paper{石橋 陽一, 宮森 恒}
{\href{https://db-event.jpn.org/deim2018/data/papers/159.pdf}{連想対話モデル : 発話文から連想した視覚情報を用いた応答文生成}}
{\upshape 第10回 データ工学と情報マネジメントに関するフォーラム (DEIM 2018)}

% ---------- Awards ----------
\cvsection{Awards and Funding}%
\entryfirst{Sponsor Award (NLP 2024)}
The 30th Annual Meeting of the Association for Natural Language Processing
\awardline{言語処理学会第30回年次大会 (NLP 2024) リクルート賞 ( 599件中 1 件 )\\ 2024}

\entry{Young Researcher Encouragement Award (NLP 2022)}
The 28th Annual Meeting of the Association for Natural Language Processing
\awardline{言語処理学会第28回年次大会 (NLP 2022) 若手奨励賞 ( 280 人中 12 人 )\\ 2022}

\entry{NAIST Research Grant (NAIST)}
\awardline{Foundation for Nara Institute of Science and Technology\\ 1,000,000 JPY for a year\\ 2021 - 2022}

\entry{Outstanding Paper Award (NLP 2020)}
The 26th Annual Meeting of the Association for Natural Language Processing
\awardline{言語処理学会第26回年次大会 (NLP 2020) 優秀賞 ( 396 件中 6 件 )\\ 2020}

\entry{Encouragement Award (YANS 2019)}
The 15th Symposium of Young Researcher Association for NLP Studies
\awardline{NLP若手の会第14回シンポジウム (YANS 2019) 奨励賞 ( 93 件中 7 件 )\\ 2019}

\entry{Excellent Research Award (IPSJ NL241)}
\awardline{Information Processing Society of Japan SIG-NL 241\\ 情報処理学会 第241回自然言語処理研究会 優秀研究賞 ( 28 件中 3 件 )\\ 2019}

\entry{President's Award}
\awardline{Kyoto Sangyo University\\ 1st in class\\ 2018}

\entry{Musubiwaza Fellowship}
\awardline{Kyoto Sangyo University\\ Tuition Waiver\\ 2015-2016, 2018}

\entry{The 20th Anniversary Special Award}
\awardline{Competition on Science and Technology Ideas (Techno-ai)\\ 2016}

\entry{Incentive Award}
\awardline{Competition on Science and Technology Ideas (Techno-ai)\\ 2016}

% ---------- Research Projects ----------
\cvsection{Research Projects}%
\entryfirst{Artificial Tactile}
\begin{indented}
Summer 2016\\[3pt]
Developing the artificial tactile device using pressure sensors and convolutional neural networks\\
\textbf{Demo video}: \url{https://youtu.be/n8lJgM3Qd-I}\\[3pt]
\textbf{Project page:} \url{https://yoichi1484.github.io/about/}\\[3pt]
\textbf{The 20th Anniversary Special Award,} Competition on science and technology ideas, 2016\\[3pt]
\textbf{Incentive Award,} Competition on science and technology ideas, 2016
\end{indented}

\end{document}
